\section{About the Committee}

\lead{The International Criminal Police Organization, known as INTERPOL, is the largest police cooperation body in the world, connecting the criminal police authorities of 196 member countries through a shared infrastructure of databases, secure communications, and international alerts.\footnote{INTERPOL. What is INTERPOL? \url{https://www.interpol.int/en/Who-we-are/What-is-INTERPOL}}}

It is also one of the most misunderstood institutions in international affairs, routinely imagined as a force of cross-border detectives able to hunt and arrest suspects anywhere on earth. The reality is narrower and, for the work of this committee, far more interesting. INTERPOL commands no officers, makes no arrests, and runs no investigations of its own\emph{.} It is a network that allows sovereign states to cooperate, and its authority ends precisely where their consent does. Grasping that distinction is the precondition for any serious debate on financial sovereignty, offshore secrecy, and the policing of illicit money.

\subsection{Mandate}

INTERPOL rests on the legal foundation laid down in its 1956 Constitution, where Article 2 gives the Organization two aims: to ensure the widest possible mutual assistance between all criminal police authorities, within the limits of national law and in the spirit of the Universal Declaration of Human Rights, and to develop the institutions that help prevent and suppress ordinary law crimes.\footnote{INTERPOL. (2024). \emph{Constitution of the ICPO-INTERPOL [I/CONS/GA/1956 (2023)]}, art. 2. \url{https://www.interpol.int/content/download/590/file/01\%20E\%20Constitution_2024.pdf}} Much of what the Organization can and cannot do is already contained in that wording. The clause “within the limits of the laws existing in the different countries” subordinates everything it does to national sovereignty, while the reference to “ordinary law crimes” deliberately leaves out the political, a boundary reinforced by Article 3, which strictly forbids any intervention of a political, military, religious, or racial character.\footnote{INTERPOL, \emph{Constitution of the ICPO-INTERPOL}, art. 3.} From these provisions flow the four principles that guide INTERPOL in practice: \textbf{national sovereignty, respect for human rights, neutrality, and constant and active cooperation.}\footnote{INTERPOL. (n.d.-b). \emph{Legal documents}. \url{https://www.interpol.int/en/Who-we-are/Legal-framework/Legal-documents}}

In practice INTERPOL is a network, not a police force. It works through three layers:

\begin{itemize}
  \item a General Assembly of all member states as its supreme body, an Executive Committee that oversees it between sessions, and a permanent General Secretariat in Lyon that runs daily operations;\footnote{INTERPOL, \emph{Constitution of the ICPO-INTERPOL}, arts. 5–32.}
  \item a National Central Bureau in each member state, the single link between that country’s police and the network, communicating through the secure I-24/7 system;\footnote{INTERPOL. (n.d.-a). \emph{Key dates}. \url{https://www.interpol.int/Who-we-are/Our-history/Key-dates}}
  \item central databases covering travel documents, fingerprints, DNA, facial images, stolen vehicles, and wanted persons, plus eight colour-coded notices, of which the Red Notice, a request to locate and provisionally detain a person pending extradition, is the best known.\footnote{Global Investigations Review. (2026). Interpol. In The practitioner’s guide to global investigations (10th ed.). \url{https://globalinvestigationsreview.com/guide/the-practitioners-guide-global-investigations/2026/article/interpol}}
\end{itemize}

Most relevant for us, the \textbf{INTERPOL Financial Crime and Anti-Corruption Centre (IFCACC)}, created in 2022, coordinates the global police response to fraud, money laundering and asset recovery, and corruption, working alongside bodies such as the Financial Action Task Force and the Egmont Group of financial intelligence units.\footnote{INTERPOL. (2022, January 25). \emph{INTERPOL launches centre against financial crime and corruption}. \url{https://www.interpol.int/en/News-and-Events/News/2022/INTERPOL-launches-centre-against-financial-crime-and-corruption}; INTERPOL. (n.d.-d). \emph{Our role in fighting financial crime}.}

\begin{keypoints}[What INTERPOL cannot do]
  \item it has no agents and no power of arrest, since every arrest is made by national police under national law;
  \item a Red Notice is not an international arrest warrant, but a request that each member may act on or decline, and many decline;\footnote{Global Investigations Review (2026), Interpol.}
  \item it does not investigate, prosecute, or determine guilt, all of which are matters for national and international courts;\footnote{INTERPOL. (2024). \emph{Repository of practice: Articles 2 and 3 of INTERPOL’s Constitution} (3rd ed.), p. 4. \url{https://www.interpol.int/content/download/12626/file/Repository\%20of\%20practice\%20Articles\%202\%20and\%203.pdf}}
  \item it is not an organ of the United Nations, but an independent intergovernmental body;
  \item its Article 3 neutrality is contested in practice, as some states have misused notices to pursue dissidents and opponents under the cover of fraud or corruption charges, forcing INTERPOL to strengthen oversight through the Commission for the Control of INTERPOL’s Files.\footnote{Parliamentary Assembly of the Council of Europe. (2017). \emph{Abusive use of the Interpol system: The need for more stringent legal safeguards} (Resolution 2161). \url{https://assembly.coe.int/nw/xml/XRef/Xref-XML2HTML-en.asp?fileid=23524&lang=en}; INTERPOL (2024), \emph{Repository of practice: Articles 2 and 3}.}
\end{keypoints}

This is the paradox at the heart of our agenda: the same instruments that let states recover stolen assets can be turned into tools of transnational repression, and a committee on financial sovereignty has to keep both in view at once.

\subsection{History}

The idea of an international police body predates INTERPOL by nearly a decade.

In 1914, the first International Criminal Police Congress met in Monaco, where officials drafted early principles for cross-border cooperation and extradition\footnote{Deflem, M. (2002). Policing world society: Historical foundations of international police cooperation. Oxford University Press.}, but the First World War suspended the project. It revived in Vienna, where on 7 September 1923 delegates from twenty countries founded the International Criminal Police Commission (ICPC) under Johannes Schober, president of the Vienna police, in answer to a new generation of criminals who used cars, telephones, and open borders to outrun investigations that still halted at the frontier.\footnote{Parliamentary Assembly of the Council of Europe. (2017). Abusive use of the Interpol system: The need for more stringent legal safeguards (Resolution 2161). \url{https://assembly.coe.int/nw/xml/XRef/Xref-XML2HTML-en.asp?fileid=23524&lang=en}} Through the interwar years, the Commission built the foundations of the modern system: a 1927 resolution urged each member to designate a central national contact point, the forerunner of today’s NCB, and by the 1930s the organization had created specialized units and launched an international radio network.\footnote{INTERPOL. (n.d.). Key dates. \url{https://www.interpol.int/en/Who-we-are/Our-history/Key-dates}}

The Second World War nearly destroyed it. After Germany annexed Austria in 1938, the Nazis seized control of the ICPC by deposing its president, and its leadership passed to a succession of senior SS officers; the headquarters were ultimately moved to Berlin, most member states withdrew, and the body was no longer functional.\footnote{Deflem (2002), The logic of nazification; INTERPOL (n.d.-a), Key dates.} \emph{This captured period is the reason the post-war organization treats political neutrality as foundational rather than decorative.} In 1946, police officials meeting in Brussels reconstituted the body as the International Criminal Police Organization, with a new headquarters in Paris, and brought “INTERPOL,” long used as its telegraphic address, into common use.\footnote{Deflem (2002), Policing world society; The Editors of Encyclopaedia Britannica. (2024). Police: International police organizations. Encyclopædia Britannica. \url{https://www.britannica.com/topic/police/International-police-organizations}} The modern Constitution followed in 1956, formally establishing the name ICPO-INTERPOL\footnote{INTERPOL, Constitution of the ICPO-INTERPOL, art. 1.} and the structure that governs the organization today, and the headquarters moved to Lyon in 1989.\footnote{The Editors of Encyclopaedia Britannica (2024), Police.}

INTERPOL’s membership and mandate have expanded steadily ever since. From twenty founding members in 1923, the organization grew to one hundred countries by 1967 and to its present 196 in 2023, making it second only to the United Nations in membership among international organizations.\footnote{Wood, M. (2026). INTERPOL as an international organization. AJIL Unbound, 120, 11–16. \url{https://doi.org/10.1017/aju.2025.10044}} Its operational focus has widened as the character of crime kept changing, which brought us to this committee agenda of transnational financial crime.
